%%Template made by Uday Khankhoje for examinations using the exam template
%%Refer to the documentation http://www-math.mit.edu/~psh/exam/examdoc.pdf 
%%for lot more bells and whistles to the standard template shown below
\documentclass[a4paper,11pt,addpoints]{exam}
\usepackage[left=1.5cm,right=1.5cm,top=1.5cm,bottom=2cm]{geometry}
%\usepackage{mathrsfs}
\usepackage{graphicx,color}
\usepackage[x11names]{xcolor}
\usepackage{epic,eepic}
\usepackage{float}
%\usepackage{mathpazo}
\usepackage{url}
\usepackage{tasks} % cria lista curta
\usepackage{multicol}
\usepackage{amsmath, amsthm, amssymb}
\pointsinmargin
\boxedpoints
\renewcommand*\half{.5}
\usepackage{setspace}
\DeclareMathOperator{\vecc}{vec}
%\renewcommand{\vec}[1]{\ensuremath{\mathbf{#1}}}
\providecommand{\sin}{} \renewcommand{\sin}{\hspace{2pt}\mathrm{sen}}
\providecommand{\tan}{} \renewcommand{\tan}{\hspace{2pt}\mathrm{tg}}

\global\vbadness=1616

\begin{document}
\noindent
%%PART 1 of header
\begin{center}
	\vspace*{-3em}
	\def\arraystretch{2.0}
	\begin{tabular}{|p{0.7\linewidth}|p{0.2\linewidth}|}
		\hline
		\textbf{Recuperação de Matemática - Primeiro Bimestre}                                                        & Pontos Obtidos $\downarrow$ \\
		\hline
		Data:\hspace{3cm}  Total de questões \textbf{\numquestions} \hspace{1cm} Total de pontos: \textbf{\numpoints} &                             \\
		\hline
		\multicolumn{2}{|l|}{Tuma: \hspace{0.3\linewidth} Nome: \hspace{0.3\linewidth} Duração: 1 hrs}                                              \\
		\hline
	\end{tabular}
\end{center}
%%PART 2 of header, if you have too many questions, this may be a problem
%%if so, use \multirowgradetable{n}[questions], where n is the number of rows you want
%%or, switch to \gradetable[h][pages] instead, 
\begin{center}
	\gradetable[h][questions]
\end{center}
%%PART 3 of header
\textbf{Instruções
	\begin{enumerate}
		\item Explique todas as questões claramente.
		\item Necessário todos os cálculos.
	\end{enumerate}
}
%%toggle comment on next line to show/hide the answers
% \printanswers
%%Now the actual paper!
\begin{questions}
	\question[2]

	Na figura abaixo, uma árvore é vista sob um ângulo de $60^\circ$, a uma distância de $60$m de sua base. Qual sua altura em metros?
	\begin{figure}[H]
		\centering
		\includegraphics[width=.4\linewidth]{figures/1.png}
	\end{figure}

	\question[2]

	Expresse em \textit{graus}:

	\begin{tasks}(2)
		\task $\dfrac{5\pi}{3}$ rad
		\task $\dfrac{8\pi}{5}$ rad
		\task $\dfrac{12\pi}{7}$ rad
		\task $\dfrac{3\pi}{10}$ rad
		\task $\dfrac{2\pi}{5}$ rad
	\end{tasks}

	\question[2]

	Expresse em \textit{radianos}:

	\begin{tasks}(2)
		\task $350^\circ$
		\task $12^\circ$
	\end{tasks}

	\question[2]

	Calcule o valor das expressões:

	\begin{tasks}(2)
		\task $E = \dfrac{\cos \dfrac{\pi}{3} + \sin \dfrac{\pi}{6}}{\tan \dfrac{\pi}{4}}$
		\task $E = \dfrac{\cos 300^\circ - \sin 210^\circ}{\tan 135^\circ}$
	\end{tasks}

	\question[2]

	Calcule o valor de:

	\begin{tasks}(2)
		\task $\cos 750^\circ$
		\task $\sin 1200^\circ$
	\end{tasks}

\end{questions}
\end{document}
